\subsection{The learned surrogate of the engine}\label{sec:m-surrogate}

\emph{Theory and source.} A learned surrogate of a phenomenological comminution model gives the same answers in one forward pass and
makes exploration of the operating space instant; Koh, Amini, McLachlan and Beaton trained a deep network to approximate the
phenomenological models of a comminution circuit for faster simulation~\cite{koh2021}. The product follows that line for one crusher.
The surrogate emulates the engine, which is itself a fast model; its purpose in the studio is a learned model that a reader can compare
with the engine, labelled as such, and the evidence of how well such a model reproduces the engine on an independent draw. It does not
emulate a plant.

\emph{Architecture.} A multilayer perceptron with GELU activations,
\begin{equation}
\hat{\mathbf y}_z = \mathrm W_4\,\mathrm{g}\!\left(\mathrm W_3\,\mathrm{g}\!\left(\mathrm W_2\,\mathrm{g}\!\left(\mathrm W_1\mathbf v + \mathbf b_1\right) + \mathbf b_2\right) + \mathbf b_3\right) + \mathbf b_4,
\qquad \mathrm{g}(a) = a\,\Phi(a),
\label{eq:mlp}
\end{equation}
with layer widths $11$-$64$-$64$-$32$-$10$ ($7{,}338$ weights and biases) and $\Phi$ the standard normal distribution function. The input
$\mathbf v$ is a one-hot block of the five machine classes followed by the six continuous inputs (setting, throw, speed, $x_{63}$, $n$ and
$A\times b$), each standardised by the mean and standard deviation of the training draw; the output is the standardised vector of
$P_{80}$, $P_{50}$, $P_{20}$, the passing at $1$, $4$, $8$, $16$ and $32$~mm, the throughput and the power, the three sizes fitted in
$\log_{10}$ and read back as ten to the power. The Bond work index of a sweep point follows the ore, $W_i=\min(20,\max(8,\ 8+0.09\,(120-A\times
b)))$~kWh/t, so the model keeps six continuous inputs.

\emph{Training.} The loss is the mean squared error in the standardised output space over the training draw $\mathcal S$,
\begin{equation}
\mathcal L(\mathrm W) = \frac{1}{|\mathcal S|\,10}\sum_{(\mathbf v,\mathbf y)\in\mathcal S}\left\|\hat{\mathbf y}_z(\mathbf v) - \mathbf y_z\right\|_2^2 ,
\label{eq:mse}
\end{equation}
minimised by Adam at a learning rate of $2\times10^{-3}$ with a weight decay of $10^{-5}$, for $500$ full-batch epochs, with torch seeded
($12345$) and single-threaded so that the run is deterministic on one machine (Algorithm~\ref{alg:surrogate}). The training points are a
Latin-hypercube sweep of the canonical engine, $1400$ per machine class over that class's own ranges (Table~\ref{tab:machines}; a gyratory
never runs a $6$~mm setting), generated by a deterministic generator (mulberry32) so that a seed yields the same points on every machine; after dropping degenerate draws $6430$ remain. The network is exported to ONNX (opset $17$, a dynamic batch axis) and
quantised to an int8 graph by dynamic quantisation of the weights.

\begin{algorithm}[!t]
\caption{Training, export and held-out evaluation of the two learned models (\texttt{stages/\allowbreak{}train.py}, \texttt{stages/\allowbreak{}evaluate.py}).}
\label{alg:surrogate}
\begin{algorithmic}[1]
\Require engine $\mathcal F$; per-class ranges; seeds $12345$ (training) and $67890$ (held out)
\State draw $1400$ Latin-hypercube points per class with seed $12345$; label each with $\mathcal F$; drop degenerate points ($6430$ remain)
\State draw $480$ points per class with seed $67890$; label; drop degenerate points ($2196$ remain)
\State standardise inputs and outputs with the training draw's statistics (float32); fit $\log_{10}$ of the three sizes
\For{epoch $=1$ \textbf{to} $500$} \Comment{surrogate, full batch}
  \State Adam step on $\mathcal L$ of Equation~\eqref{eq:mse} (rate $2\times10^{-3}$, weight decay $10^{-5}$)
\EndFor
\For{epoch $=1$ \textbf{to} $400$} \Comment{autoencoder, full batch}
  \State $\tilde{\mathbf h}\gets\mathbf h+\boldsymbol\epsilon$, $\boldsymbol\epsilon\sim N(0,0.05^2\mathrm I)$; Adam step on $\|\mathrm{AE}(\tilde{\mathbf h})-\mathbf h\|^2$
\EndFor
\State threshold $\gets$ 99th percentile of the training reconstruction scores
\State export both to ONNX (opset 17); quantise the surrogate to int8; record the export parity and the SHA-256 of each graph
\State evaluate the shipped graphs through onnxruntime on the held-out draw: $R^2$ and MAPE per output; false-alarm and detection rates
\end{algorithmic}
\end{algorithm}

\emph{Held-out evaluation.} The held-out set is an independent second Latin-hypercube draw (seed $67890$, $480$ points per class, $2196$
after dropping degenerate draws), not a split of the training draw. Every number is computed through onnxruntime on the shipped float32 graph
against the canonical engine, and the MAPE also on the int8 graph, so the metrics describe what the browser runs. Monotonicity is checked on
fifteen curves ($16$ settings across each class's range, at the $25$, $50$ and $75\%$ points of its other inputs): $P_{80}$ must not fall by more
than $0.001$~mm as the setting opens. The export parity is recorded twice: the trained torch module against the float32 graph on the held-out
draw, and the float32 graph against the int8 graph. The results are in Section~\ref{sec:results-learned}.

\emph{Cost.} A forward pass is about $7.3\times10^{3}$ multiply-adds; training both models, with the two sweeps and the evaluation, took $\resModelsS$~s on
the workstation's processor (Table~\ref{tab:resources}). \emph{Applicability.} The slot answers only an operating point that is not calibrated and lies inside the training
ranges; otherwise the engine answers, labelled \emph{calibrated case} or \emph{outside the training range}. The surrogate returns one value per
output and no uncertainty.

\emph{Implementation and tests.} \texttt{model/\allowbreak{}surrogate.py}, \texttt{model/\allowbreak{}sweep.py} (\texttt{RANGES}, \texttt{draw}), \texttt{stages/\allowbreak{}train.py},
\texttt{stages/\allowbreak{}evaluate.py}, \texttt{model/\allowbreak{}bundles.py} (the release rule: the slot is released only when its contract validates and every graph
exists with its SHA-256 and size), and \texttt{model/\allowbreak{}crusher/\allowbreak{}encoding.py}, which reproduces the browser's encoding byte for byte. The acceptance
criteria are an $R^2$ for every output computed from the shipped graph, $P_{80}$ monotone in the setting on every class, and the engine answering,
labelled, when the bundle is invalid (\texttt{test\_metrics\_validate} and a test that swaps the slot to its engine baseline).

\subsection{The denoising autoencoder as an extrapolation score}\label{sec:m-ae}

\emph{Theory.} An autoencoder trained to reconstruct the signature of engine results reconstructs badly a result that lies far from the manifold
it was trained on; its reconstruction error is then an anomaly score that says the learned surrogate is extrapolating. Training it on inputs
corrupted by noise and asking for the clean input (a denoising autoencoder) makes the reconstruction a projection onto the training manifold rather
than an identity map. A product of this kind is motivated by the use of the power-draw signal of a cone crusher to diagnose feed segregation, a
real anomaly of the machine~\cite{grondahl2018}; the product's score detects distance from the engine's sampled envelope, not a mechanical fault.

\emph{Architecture and score.} The encoder maps the fourteen standardised features $\mathbf h$ (the $\log_{10}$ of $P_{80}$, $P_{50}$ and $P_{20}$,
the passing at five sizes, throughput, power, reduction ratio, $\log_{10}F_{80}$, specific energy and the feed exponent) through widths $14$-$16$-$6$
with GELU, and the decoder maps back through $6$-$16$-$14$ ($692$ weights and biases). The score and the alarm are
\begin{equation}
a(\mathbf h) = \frac{1}{14}\left\|\mathrm{AE}(\mathbf h) - \mathbf h\right\|_2^2,\qquad
\text{alarm} \iff a(\mathbf h) > a_{99}\ \lor\ a(\mathbf h)\notin\mathbb R,
\label{eq:ae}
\end{equation}
with $a_{99}=0.24454$ the $99$th percentile of the scores of the training draw (whose mean score is $0.05181$). Training uses the same $6430$ points
as the surrogate, Adam at $2\times10^{-3}$ with weight decay $10^{-5}$ and $400$ full-batch epochs (Algorithm~\ref{alg:surrogate}), on the loss
\begin{equation}
\mathcal L_{\mathrm{AE}} = \frac{1}{14\,|\mathcal S|}\sum_{\mathbf h\in\mathcal S}\left\|\mathrm{AE}(\mathbf h+\boldsymbol\epsilon)-\mathbf h\right\|_2^2,
\qquad \boldsymbol\epsilon\sim N\!\left(\mathbf 0,\ 0.05^2\,\mathrm I_{14}\right),
\label{eq:dae}
\end{equation}
with a fresh noise draw at every epoch, so that the network learns to map a perturbed signature back to the training manifold. The noise of
$0.05$ is small against the unit standard deviation of each standardised feature, so the reconstruction of a training point stays close
to the point (the mean training score of $0.052$ is about a twentieth of the variance of a feature), while a point that no training signature
resembles is mapped to the nearest one and scores high.

\emph{Why the features are engine outputs.} The fourteen features are results of the engine, not its inputs: the score asks whether the
product sizes, the readouts, the throughput, the power, the reduction ratio and the specific energy of a point are jointly typical of the
training draw. A point whose inputs are each inside their training ranges but whose combination gives an unusual product, for example a fine
feed at a wide setting, can therefore raise an alarm, while a point whose inputs lie outside the ranges but whose results resemble the training
results does not. This is the complement of the envelope check, which tests inputs one at a time against rules.

\emph{Evaluation protocol.} The false-alarm rate is the share of the $2196$ held-out points that raise an alarm. The detection rate is measured on
$2196$ planted points, one per held-out point, whose operating inputs are moved outside that class's training range and then run through the
engine: by the held-out index modulo $18$, one of the six continuous inputs is moved above its upper bound by half the range width (patterns $0$ to
$5$), or below its lower bound to half the bound ($6$ to $11$), or two neighbouring inputs are moved, the first above and the second below ($12$ to
$17$); a planted point whose engine result has a non-finite feature is left out. The envelope check is scored on the same points; on the synthetic
lane it flags only implausible inputs, so it is the baseline the score must complement.

\emph{Parameters and how they were set.} The bottleneck width $6$ and the noise $0.05$ are the product's choices; the threshold is set by the
training scores, so the expected false-alarm rate on new data from the training distribution is about one percent. \emph{Cost.} About
$6.6\times10^{2}$ multiply-adds per score. \emph{Limits.} The training data are engine outputs; the one percent is measured on one held-out draw and
is not a guaranteed rate on other data.

\emph{Implementation and tests.} \texttt{model/\allowbreak{}psd\_ae.py}, \texttt{stages/\allowbreak{}train.py}, \texttt{stages/\allowbreak{}evaluate.py} and
\texttt{model/\allowbreak{}bundles.py} (\texttt{anomaly\_score}, \texttt{reconstruction\_score}, \texttt{applicability}). The acceptance criteria are the
false-alarm and detection rates on the held-out and planted points reported beside those of the envelope check, and the envelope check answering,
labelled, when the bundle is invalid (\texttt{test\_metrics\_report\_the\_anomaly\_comparison\_and\_the\_planted\_points} and a test that
swaps the slot to the envelope check).

\subsection{The drift projection: hours until the product leaves its size target}\label{sec:m-drift}

\emph{Theory and source.} As the liners of a cone wear, the closed-side setting opens and the product coarsens; plants compensate by re-setting
the machine. One sourced cone-liner wear model is the shear-augmented Archard form of Lindqvist and Evertsson~\cite{lindqvist2006}, extended
with a time term by Ma, Fan and He~\cite{ma2016}, in which the wear of a crushing zone after $t$ hours is the wear near the setting point scaled by
the zone's relative stress and growing with the cumulative tonnage; in the plant of that study the setting was re-set every $45$~h. The parameters
of that model, a liner wear rate in grams per tonne, a liner life in hours or a drift of the setting in millimetres per hour were not found in any
source the product could read, so the product models no wear and asserts no drift rate. What it computes is the opening of the setting that a
product size target allows, from the engine's own $P_{80}$, and the hours that opening lasts at a drift rate the user supplies:
\begin{equation}
s^{\ast} = \min\left\{ s' \ge s_0 : P_{80}(\mathbf u; s') \ge P_{80}^{\mathrm{target}} \right\},\qquad
\Delta s = s^{\ast} - s_0,
\label{eq:target}
\end{equation}
\begin{equation}
t_{out} = \frac{\max(\Delta s,\ 0)}{r},\qquad r > 0 ,
\label{eq:hours}
\end{equation}
with $s_0$ the current setting, $P_{80}(\mathbf u;s')$ the engine's product size at the operating point with the setting replaced by $s'$, and
$r$ the drift rate in mm/h. The view plots $t_{out}$ against $r$ on logarithmic axes from $0.01$ to $1$~mm/h, with the target as a control that opens
at the case's $P_{80}$ plus $10\%$ (a control value, not a claim) and the $45$~h interval of that plant as the one reference line
(Figure~\ref{fig:drift} in Section~\ref{sec:results-drift}).

\emph{Numerical scheme.} Algorithm~\ref{alg:target} finds $s^{\ast}$ by a bracket and a bisection. It samples the engine on $64$ settings from $s_0$
to a ceiling $s_{\max}$, stops at the first setting the engine marks invalid (for example a pass-through feed), requires that $P_{80}$ never falls
over the valid samples by more than four units in the last place, and bisects $60$ times between $s_0$ and the last valid sample. Sixty halvings of a
bracket of at most $240$~mm leave an interval below $10^{-15}$~mm, so the result is limited by floating-point resolution and both engines run the
same iterations. The ceiling is the upper bound of the class's training range: $90$~mm for the secondary cone, $22$ tertiary, $16$ short-head, $160$
jaw and $240$ gyratory.

\begin{algorithm}[!t]
\caption{The setting at which the engine's $P_{80}$ reaches a target (\texttt{css\_at\_p80\_target}).}
\label{alg:target}
\begin{algorithmic}[1]
\Require operating point $\mathbf u$ with setting $s_0$; target $P^{\ast}>0$; ceiling $s_{\max}>s_0$
\Ensure $s^{\ast}$, or no answer
\State $p_0\gets P_{80}(\mathbf u;s_0)$; \textbf{if} $\mathbf u$ is invalid \textbf{then return} no answer
\If{$p_0\ge P^{\ast}$} \Return $s_0$ \Comment{target already reached: zero hours}
\EndIf
\State $(s_{\mathrm{last}}, p_{\mathrm{last}})\gets(s_0,p_0)$
\For{$i=1$ \textbf{to} $63$}
  \State $s\gets s_0+(s_{\max}-s_0)\,i/63$; \textbf{if} the engine marks $s$ invalid \textbf{then break}
  \State \textbf{if} $P_{80}(\mathbf u;s)<p_{\mathrm{last}}-4\,\mathrm{ulp}(p_{\mathrm{last}})$ \textbf{then return} no answer \Comment{not monotone}
  \State $(s_{\mathrm{last}}, p_{\mathrm{last}})\gets(s, P_{80}(\mathbf u;s))$
\EndFor
\State \textbf{if} $p_{\mathrm{last}}<P^{\ast}$ \textbf{then return} no answer \Comment{never reached at a valid setting}
\State $(l,h)\gets(s_0,s_{\mathrm{last}})$
\For{$k=1$ \textbf{to} $60$}
  \State $m\gets(l+h)/2$; \textbf{if} $m$ is invalid or $P_{80}(\mathbf u;m)\ge P^{\ast}$ \textbf{then} $h\gets m$ \textbf{else} $l\gets m$
\EndFor
\State \Return $(l+h)/2$
\end{algorithmic}
\end{algorithm}

\emph{The shared table.} The $45$ rows of the reference table hold, for each of nine operating points, three targets above the current
$P_{80}$ ($+3$, $+10$ and $+30\%$), one target already exceeded and one that is never reached. For S01 they give the settings $33.043$, $35.212$ and
$40.673$~mm for the first three, the current $32$~mm for the exceeded target (zero hours), and no answer for a target of $10^{6}$~mm, which no valid
setting below the $90$~mm ceiling reaches. Both engines must reproduce every row to $10^{-6}$~mm, and a further test, which does not use the table,
evaluates the engine at each returned setting and requires its $P_{80}$ to equal the target. A patched engine whose $P_{80}$ falls over the bracket
must return no answer, so the monotonicity check is itself tested.

\emph{Parameters.} None with a default for the drift: no function and no control of the view has a default rate. The $64$ samples, the $60$
halvings and the tolerance of four units in the last place are the engine's. \emph{Cost.} At most $64+60$ engine evaluations, about
$\resDrift$~s at the median engine time. \emph{Limits.} Not a wear model and not a set point; the projection is the engine's (or the calibrated engine's)
$P_{80}$ at the opened setting, with the limits of the engine, and it ignores the profile distortion of worn liners and their work hardening.

\emph{Implementation and tests.} \texttt{css\_at\_p80\_target} and \texttt{hours\_to\_target} in \texttt{model/\allowbreak{}crusher/}, and their
TypeScript port. The acceptance criteria are that both engines return the same setting on a shared table of $45$ rows (nine operating points,
including the calibrated lane, at five targets each) to $10^{-6}$~mm, that the engine's $P_{80}$ at the returned setting is the target, and that a
target already exceeded, never reached, an invalid point and a non-monotone $P_{80}$ give zero hours or no answer, never a guess
(\texttt{test\_the\_shared\_table\_is\_reproduced\_to\_1e6\_mm}, \texttt{test\_a\_target\_never\_reached\_below\_the\_ceiling\_is\_none}).

\subsection{CrusherBayes: Bayesian calibration with a structural discrepancy}\label{sec:m-crusherbayes}

\emph{Idea.} The calibration of Rocha et al.\ is a point fit~\cite{rocha2024}. CrusherBayes, the method the product proposes, places a prior on the
engine's fitted scalars, adds an explicit model-form discrepancy per target in the form of Kennedy and O'Hagan~\cite{kennedy2001}, and returns a
posterior predictive distribution for a survey that was not used in the fit. The prediction stated before the computation was that this
distribution would be better calibrated and sharper than the residual band of CrusherCal, and would show which scalars the ten surveys identify and
which stay at their prior. The method is derived and analysed in its own manuscript~\cite{crusherbayes2026}; this subsection gives what a reader of
the product needs to reproduce it. The searches that preceded it found no Bayesian calibration with discrepancy of a cone-crusher
population-balance model evaluated held out; they were not exhaustive (Section~\ref{sec:limits}).

\emph{Model.} For survey $i$ and target $t$ (throughput in t/h, power in kW; $P_{80}$ is not measured and is never a target),
\begin{equation}
y_{t,i} = f_t\!\left(\boldsymbol\vartheta;\ s_i, f_{80,i}\right) + \delta_t(\mathbf z_i) + \varepsilon_{t,i},\qquad
\delta_t(\mathbf z) = d_{t0} + d_{t1}\,z_{\mathrm{CSS}} + d_{t2}\,z_{f_{80}},
\label{eq:cbmodel}
\end{equation}
\begin{equation}
d_{tj}\sim N(0,\tau_t^2),\qquad \varepsilon_{t,i}\sim N(0,\sigma_t^2),\qquad \tau_Q = 140\ \mathrm{t/h},\quad \tau_P = 36\ \mathrm{kW},
\label{eq:cbprior}
\end{equation}
where $f_t$ is the engine at the eight scalars $\boldsymbol\vartheta=(Q_T, W_i, \zeta, P_n, s_1, s_2, K_3, \Delta t_{10})$ of
Table~\ref{tab:priors}, with $s_1$ and $s_2$ multipliers of $K_1$ and $K_2$ and $\Delta t_{10}$ a shift of the constant $64$ of the $t_{10}$
relation, and $\mathbf z$ the setting and $f_{80}$ standardised on the training surveys of the fold. The discrepancy scales are about $12\%$ of
the typical level of each target. Given $\boldsymbol\vartheta$ the discrepancy and the noise are Gaussian, so the discrepancy is integrated out
exactly: with $\mathrm Z$ the $9\times3$ design matrix $[\mathbf 1, \mathbf z_{\mathrm{CSS}}, \mathbf z_{f_{80}}]$ of the training surveys,
\begin{equation}
\mathbf y_t \mid \boldsymbol\vartheta,\sigma_t \ \sim\ N\!\left(f_t(\boldsymbol\vartheta),\ \boldsymbol\Sigma_t\right),\qquad
\boldsymbol\Sigma_t = \sigma_t^2\,\mathrm I + \tau_t^2\,\mathrm Z\mathrm Z^{\top},
\label{eq:cbmarginal}
\end{equation}
and the sampler explores ten dimensions, the eight scalars and the two log noise scales. The log posterior is the sum of the two Gaussian log
densities of Equation~\eqref{eq:cbmarginal}, computed by a Cholesky factorisation of $\boldsymbol\Sigma_t$, and the truncated-normal log prior. The
discrepancy coefficients are drawn afterwards from their exact Gaussian conditional,
\begin{equation}
\mathbf d_t\mid\boldsymbol\vartheta,\sigma_t,\mathbf y_t\ \sim\
N\!\left(\tau_t^2\mathrm Z^{\top}\boldsymbol\Sigma_t^{-1}\mathbf r_t,\ \ \tau_t^2\mathrm I - \tau_t^4\,\mathrm Z^{\top}\boldsymbol\Sigma_t^{-1}\mathrm Z\right),\qquad
\mathbf r_t = \mathbf y_t - f_t(\boldsymbol\vartheta),
\label{eq:cbcond}
\end{equation}
and the predictive distribution of a held-out survey with design row $\mathbf z_h=[1,\ z_{\mathrm{CSS},h},\ z_{f_{80},h}]$ (the
leading one carries the discrepancy constant) at one posterior sample is Gaussian with mean
$f_t(\boldsymbol\vartheta;s_h,f_{80,h})+\mathbf z_h^{\top}\mathbb E[\mathbf d_t]$ and variance $\sigma_t^2+\mathbf z_h^{\top}\mathrm{Cov}[\mathbf
d_t]\,\mathbf z_h$. The engine is its own simulator: at a median of $\resEngineMed$~ms per evaluation no emulator is needed.

\emph{Priors.} Each engine scalar has a truncated normal prior centred on the engine's own value, on the log scale where marked
(Table~\ref{tab:priors}): the nominal values are those of the source calibration and of the itabirite blend, and the widths are weakly informative
choices of the author that bracket the quoted spans, because no published range of these fitted scalars for this ore was found. The noise has
$\ln\sigma_Q\sim N(\ln 60,1)$ and $\ln\sigma_P\sim N(\ln 15,1)$, about $5\%$ of the typical level, truncated to $[\ln1,\ln1000]$.

\emph{Sampler.} The ensemble sampler of \texttt{emcee} $3.1.6$~\cite{emcee2013} with the differential-evolution moves of ter
Braak~\cite{terbraak2006} (weight $0.8$) and the snooker update of ter Braak and Vrugt~\cite{terbraak2008} (weight $0.2$): gradient-free, as the
engine requires. The affine-invariant stretch move of Goodman and Weare~\cite{goodman2010} was tried first and did not converge: its autocorrelation
time stayed above $100$ steps and grew with the chain, along the parameters the surveys do not identify. Algorithm~\ref{alg:crusherbayes} gives one
posterior. Twenty walkers start at the prior centre perturbed by $0.05$ prior standard deviations, run $7000$ steps, discard $1000$, and are thinned
by $10$, which leaves $12{,}000$ samples. A posterior is recorded as converged when the post-burn chain is at least $50$ integrated autocorrelation
times long, the largest over the parameters; a posterior that fails is recorded as not converged and is never kept silently. The $80\%$ interval of
a held-out survey is the $10$th to $90$th percentile of $20$ predictive draws per retained sample.

\begin{algorithm}[!t]
\caption{One CrusherBayes posterior and its held-out prediction (\texttt{hp500/\allowbreak{}crusherbayes.py}, \texttt{\_fit}).}
\label{alg:crusherbayes}
\begin{algorithmic}[1]
\Require training surveys $\mathcal T_{-h}$ (or all ten), held-out survey $h$ (or none), seed; priors of Table~\ref{tab:priors}
\State standardise the setting and $f_{80}$ on $\mathcal T_{-h}$; build $\mathrm Z$
\State $\log\pi(\boldsymbol\theta)\gets\log p(\boldsymbol\theta)+\sum_{t\in\{Q,P\}}\log N\!\left(\mathbf y_t;\ f_t(\boldsymbol\vartheta),\ \sigma_t^2\mathrm I+\tau_t^2\mathrm Z\mathrm Z^{\top}\right)$ \Comment{engine on every training survey}
\State initialise $20$ walkers at the prior centre plus $0.05$ prior standard deviations of Gaussian noise
\State run $7000$ steps of the ensemble sampler with DE (weight $0.8$) and DE-snooker (weight $0.2$) moves
\State $\hat\tau_{\max}\gets$ largest integrated autocorrelation time after $1000$ steps; converged $\gets 6000\ge50\,\hat\tau_{\max}$
\State retain the chain after $1000$ steps, thinned by $10$ ($12{,}000$ samples)
\If{$h$ is given}
  \For{each retained sample}
    \State compute the exact conditional of $\mathbf d_t$ (Equation~\eqref{eq:cbcond}); predictive mean and variance at survey $h$
    \State draw $20$ predictive values
  \EndFor
  \State \Return mean of the predictive means; $10$th and $90$th percentiles of the draws; log predictive density
\Else
  \State draw $\mathbf d_t$ once per sample from Equation~\eqref{eq:cbcond}; \Return the samples for the contraction table
\EndIf
\end{algorithmic}
\end{algorithm}

\emph{Evaluation and adoption rule.} Strict leave-one-survey-out as in CrusherCal: one posterior per fold from nine surveys, the held-out survey only
in the prediction, and one more posterior from all ten for the contraction table, the ratio of posterior to prior standard deviation of every
parameter. Held out, per target: the MAPE of the predictive mean, the RMSE, the mean log predictive density, the coverage of the $80\%$ interval and
its mean width, beside the setting-only line, the calibrated backbone and the free-form fit scored with their nested bands. The adoption rule,
written before the run and applied by code, drops the method as a predictor if (a) its power MAPE is not below the line's by at least three points
and its throughput MAPE does not beat the backbone's by at least three points, and (b) on either target its coverage lies outside $0.7$ to $0.9$ or
its mean band is wider than the backbone's band. The contraction table is reported in either case. The results are in
Section~\ref{sec:results-cb}.

\emph{Cost.} Each log-posterior evaluation runs the engine on the nine (or ten) training surveys, so one posterior costs $140{,}000\times9$ engine
evaluations; the bake spreads the walkers of the eleven posteriors over the processor's threads, and the whole step takes about
$\resHpMin$~minutes on the workstation (Table~\ref{tab:resources}). \emph{Limits.} With ten surveys and two targets the
posterior of most scalars is the prior; the discrepancy can absorb what the engine cannot express, which is the known confounding of calibration
and discrepancy; and the power form of the engine is the weak part, so most of any gain on power may come from the discrepancy, which the
setting-only line already approximates.

\emph{Implementation and tests.} \texttt{hp500/\allowbreak{}crusherbayes.py} (\texttt{PRIORS}, \texttt{log\_prior}, \texttt{\_Problem.log\_prob},
\texttt{\_Problem.conditional}, \texttt{\_Problem.predict}, \texttt{\_fit}, \texttt{kill\_criterion}); the engine takes the scalars through one optional
argument whose defaults are the shipped values, so the TypeScript engine and the shared reference are untouched. The tests check that the default
scalars are the shipped engine, that each prior is centred on the engine's value with a source, the prior bounds, that the folds never see their
held-out survey, the seed, that the document is byte-identical and independent of the number of worker processes, that a short chain is recorded as
not converged, the discrepancy marginal against a Monte Carlo of the model, and the adoption rule in every combination of its terms.

\begin{table}[!t]
\centering
\caption{Priors of the eight CrusherBayes engine scalars: truncated normal around the engine's own value, on the scale shown, with the truncation on
that scale. Nominal values are those of the source calibration and the itabirite blend; widths are the author's. Generated from
\texttt{crusherbayes.json}.}
\label{tab:priors}
\footnotesize
% generated by figures/make_tables.py; do not edit by hand
\begin{tabular}{@{}l l r l r l@{}}
\toprule
Scalar & Unit & Nominal & Scale & Prior sd & Truncation (on the scale) \\
\midrule
$Q_T$, capacity base & t/h & $647.00$ & log & $0.25$ & $[5.704,\ 7.244]$ \\
$W_i$, blend work index & kWh/t & $9.00$ & log & $0.17$ & $[1.609,\ 2.639]$ \\
$\zeta$, power slope & -- & $1.30$ & log & $0.23$ & $[-0.693,\ 0.916]$ \\
$P_n$, no-load power & kW & $110.00$ & linear & $50.00$ & $[0.000,\ 300.000]$ \\
$s_1$, multiplier of $K_1$ & -- & $1.00$ & log & $0.15$ & $[-0.693,\ 0.693]$ \\
$s_2$, multiplier of $K_2$ & -- & $1.00$ & log & $0.15$ & $[-0.693,\ 0.693]$ \\
$K_3$, classification exponent & -- & $2.30$ & linear & $0.50$ & $[1.000,\ 4.000]$ \\
$\Delta t_{10}$, shift of the constant 64 & points & $0.00$ & linear & $4.00$ & $[-15.000,\ 15.000]$ \\
\bottomrule
\end{tabular}

\end{table}
